\section{Evaluation Methodology}
\label{sec:method}

The experiments need no language model. The current-code results were produced on 2026-10-04 on a Linux container with a 4-core Intel Xeon at 2.1\,GHz, 15\,GiB RAM and no GPU, with Python~3.11.15, Playwright~1.56.0, headless Chromium~141 and LangGraph~1.2.12; the decisions of the code before the fixes come from an earlier run on a similar container and are deterministic. Scripts, seeds and raw outputs are under \texttt{artifacts/} and \texttt{benchmarks/e2e/} in the repository.

\subsection{Gate Study}
\label{sec:method-gate}

\paragraph{States and splits} A state is a goal, a URL, an HTML document served through Playwright request interception, the agent's answer and plan state where relevant, and an oracle label. Every state ends with a dispatched action, so a dispatch-only rule accepts all of them. The \emph{design} set (38 states in six categories) was used to find the gate's defects and design the fixes. \emph{Held-out~1} (12 states) was written after the fixes and not used to change them. \emph{Held-out~2} (42 states, 22 unsatisfied) was written to a fixed procedure---varied user phrasings per category, including wordings the gate's trigger lists do not contain, with a balanced split---and committed before the gate was run on it; no gate code changed afterwards. Its SHA-256 is recorded in the repository.

\paragraph{Oracle} The label $\Phi^{\star}$ of a state states whether the goal is achieved by the page and answer shown, judged from the page's facts (for example, whether the answer contains the year in the article). It was written together with the state, before the state was judged by the gate, and it is computed by no code shared with $\Phi$. Its limitation is that one author wrote all states and labels, so no inter-annotator agreement exists; a protocol for independent relabelling with Cohen's $\kappa$~\cite{cohen1960coefficient} is in the repository.

\paragraph{Modes and statistics} Each state is judged by the pre-check and $\Phi$ called directly, and by the real verification node (database, browser-pool and evidence I/O stubbed) with an empty or a generic one-sentence planner reasoning. For the ablation, each conjunct of $\Phi$ and the pre-check $\beta$ is removed in turn by a configuration switch. Counts carry Wilson and Clopper--Pearson 95\% intervals~\cite{wilson1927probable,clopper1934use}; comparisons on the same states use the exact McNemar test~\cite{mcnemar1947note} on discordant pairs. Latency is the median of 20 repetitions after a warm-up.

\subsection{End-to-End Harness}
\label{sec:method-e2e}

\paragraph{Environment and tasks} Live websites change and block automated traffic, and WebArena-style site images were not available offline, so we built five self-hosted sites served on the loopback interface with fictional content that a model cannot answer from memory: a search engine, an encyclopedia, a shop with search and cart, a town-office contact form and newsletter, and a notes application; a further path redirects to a sign-in wall. The 40 tasks (10 development, 30 test) cover search, navigation, text entry, extraction, form submission, multi-step tasks and 5 unsatisfiable requests, at four difficulty tiers by the number of actions a solution needs (1, 2--3, 4--6, more). The notes site appears only in the test split. Tasks needing an account, payment, a CAPTCHA, the live web, or an outcome that neither server state nor answer text can verify were excluded by design.

\paragraph{Oracle} $\Phi^{\star}$ is computed from signals that neither the agent nor $\Phi$ uses: server-side records of the sites (cart contents, form submissions, saved notes, field values autosaved by page scripts), the final URL, and the ground-truth facts in the site data matched against the agent's answer; unsatisfiable tasks are never achieved. It calls no model and shares no code with the verifier. We measured its error rate on 114 scripted trajectories driven directly through Playwright without the agent: a correct solution for each task, an idle trajectory that only opens the start page, and a near miss for 34 tasks (wrong value, wrong item, extra submission, partial query).

\paragraph{Protocol} Table~\ref{tab:e2e} (Appendix~\ref{app:tasks}) lists the configurations: the full system, two baselines---(a) no completion check and (b) self-verification by the planner model---and one-factor ablations of each conjunct, memory, recovery and each recovery strategy. Each configuration runs in its own process with its own database and memory store, at temperature 0 with a fixed seed, three trials per task, and the same budget of 15 planner calls. Per episode the harness records the oracle verdict, the agent's status, browser actions, model calls and tokens, wall-clock time, each verification decision with its latency, recovery strategies, retrieved memories, Ollama version and model digests, and every host the browser contacted. The analysis reports success, FCR and FRR with Wilson intervals, per-category breakdowns, exact McNemar tests and task-level bootstrap intervals~\cite{efron1979bootstrap} for differences to the full system, and the overhead of $\Phi$ relative to baseline (a). A published agent framework as a third baseline would need its own model and prompts and is left to the authors' runs.


\subsection{Simulation-Based Evaluation and Data Provenance}
\label{sec:sim}

\paragraph{Why simulation} The end-to-end protocol needs a local language model, which could not be installed in the environment used for this paper. To study the questions of Section~\ref{sec:hypotheses} before those runs exist, we built a seeded simulation of the agent's completion decision, recovery, memory and latency (\texttt{benchmarks/simulation}). The simulation does not execute the agent. Its outputs follow from two kinds of input, and every number in this paper carries one of five classes: \textsc{actual} (measured on the implementation), \textsc{derived} (computed from actual measurements), \textsc{synthetic} (produced by the simulation), \textsc{literature}, or \textsc{proposed} (behaviour of an unimplemented component). Tables and figures with synthetic values say so in their captions, and an evidence ledger lists every simulated value with its source.

\paragraph{Measured inputs} Without a model we measured, on the implementation: the decisions of $\Phi$ and of a basic success detector (error check $E$ only) on the 92 gate states; the latency of navigation, clicks, form filling, DOM extraction, screenshots and $\Phi$ on the self-hosted sites (30 repetitions each); the failure class and strategy sequence that the real recovery engine assigns to six failure messages; and the ranks at which the real Chroma memory store returns relevant records for ten queries over a 23-record synthetic memory corpus. Detection rates used by the simulation are derived from the two held-out gate sets only, with the estimator $(k+0.5)/(n+1)$.

\paragraph{Assumptions and generation} What needs a model---how often the planner claims completion of an unfinished task, how often actions fail and are recovered, how much memory helps---is assumed and listed with its class in Table~\ref{tab:sim-params} (Appendix~\ref{app:simparams}). A scenario has a task category (navigation, search, form submission, data entry, extraction, multi-step) and an outcome class (successful, incomplete, incorrect action, missing evidence, contradictory evidence, stale evidence, ambiguous) that sets the probability that the goal is actually met. Each scenario draws its random numbers once, and every configuration reuses them, so comparisons are paired by scenario. The true outcome is evaluated per configuration because memory and continuation after a rejection change the final state. Parameters and code were committed (\texttt{3a0aa8e}) before the reported held-out set (seed 202, 100 scenarios) was generated; a development set (seed 101) was used only to debug the code. Monte Carlo replications regenerate the scenario set 1{,}000 times to show the variability the generator itself induces.

\paragraph{Metrics} With TP, FP, FN, TN the confusion counts of reported completion against the true outcome over $N$ scenarios: $\mathrm{TSR}=\mathrm{TP}/N$; FCR and FRR as in Eq.~\eqref{eq:fcr}, so that precision $=1-\mathrm{FCR}$ and recall $=1-\mathrm{FRR}$; F1 and accuracy $(\mathrm{TP}+\mathrm{TN})/N$. Recovery success is recovered over recoverable failures; retrieval quality is precision and recall at $k=3$ and the mean reciprocal rank (MRR) of the first relevant record. Episode time is decomposed as
\begin{equation}
T = T_{\mathrm{inf}} + T_{\mathrm{browser}} + T_{\mathrm{evidence}} + T_{\mathrm{verify}} + T_{\mathrm{recovery}},
\label{eq:time}
\end{equation}
with browser, evidence and $\Phi$ times resampled from the measurements and inference time assumed. Paired configurations are compared by the exact McNemar test on per-scenario indicators, differences by a scenario-level bootstrap interval, and success rates by Cohen's $h$. A $p$-value here describes the synthetic sample under the stated assumptions; it is not evidence about the deployed agent.

\paragraph{Independent validation} A separate validation agent recomputed all confusion counts, rates, paired tests, memory and recovery statistics from the per-scenario files with independently written code and found no discrepancy (\texttt{simulation\_validation\_report.md}).
